\documentclass{article}
\usepackage[utf8]{inputenc}
\usepackage[T1]{fontenc}
\usepackage{geometry}
\geometry{
    a4paper,
    left=2.5cm,
    right=2.5cm,
    top=2.5cm,
    bottom=2.5cm
}
\usepackage{graphicx}
\usepackage{float}
\usepackage{hyperref}
\usepackage[english]{babel}
\usepackage{amsmath}
\usepackage{amssymb}
\usepackage{array}
\usepackage{booktabs}
\usepackage{caption}
\usepackage{subcaption}
\usepackage{enumitem}
\usepackage{fancyhdr}
\usepackage{titlesec}
\usepackage{url}
\usepackage{textcomp}
\usepackage{xcolor}
\usepackage{listings}
\usepackage{multirow}
\usepackage{longtable}
\usepackage{tabularx}
\usepackage{times}

% Set up page style
\pagestyle{fancy}
\fancyhf{}
\rfoot{\thepage}
\renewcommand{\headrulewidth}{0pt}

% Title formatting
\titleformat{\chapter}[display]{\normalfont\bfseries}{\chaptertitle\thechapter}{0pt}{\Huge}
\titlespacing*{\chapter}{0pt}{50pt}{40pt}

% Document info
\title{Template Document}
\author{Academic Support Agent}
\date{\today}

\begin{document}
\maketitle

\fontfamily{cmss}
\fontsize{12}{18}
\selectfont

\begin{center}
	\includegraphics[scale=0.18]{brasao2_vertical}
\end{center}

\begin{center}
	\vspace{14ex}
	Campus de Crate s\\	
	Est gio Supervisionado\\
	\vspace{10ex}
	\textbf{{
			\large PLANO DE TRABALHO DE EST GIO SUPERVISIONADO
	}}\\
	\vspace{16ex}
	{\large 
		Atividade de Est gio Supervisionado\\
		Est gio <Especificar Institui o>\\
		<T tulo do Projeto que est  sendo desenvolvido>
	}
\end{center}

\newpage

\section*{Identifica o}

\textbf{Estagi rio:}\\

Nome:  \\
Matr cula:  \\
Telefones:  \\
E-mail:  \\
Cargo:  \\

\textbf{Institui o do Est gio:}\\

Nome:  \\
Endere o:  \\
Ramo de Atividade:  \\
Supervisor:  \\
Telefone:  \\
E-mail:  \\

\textbf{Professor Orientador:}

\newpage



\newpage

\addcontentsline{toc}{section}{Introdu o}
\section*{Introdu o}

\sethlcolor{lightgray}

Este documento tem por objetivo mostrar o trabalho desenvolvido pelo aluno(a) \hl{<Nome do aluno(a)>} durante o est gio no(a) \hl{<Nome da empresa>}.

\addcontentsline{toc}{section}{Empresa}
\section*{Empresa}

\hl{<Adicionar suas impress es em rela o   empresa. Pontos positivos e negativos do local de est gio> OBS.: N o   aquele texto padr o da empresa, s o suas impress es. \\
<Relato sobre processos e ferramentas adotados pelas empresas.>
}

\addcontentsline{toc}{section}{Relato de Experi ncia}
\section*{Relato de Experi ncia}

\hl{<Fazer um relato de experi ncia descrevendo aqui tudo como foi seu est gio.>}

\subsection*{Cronograma de atividades}

\hl{<As atividades descritas no relato de experi ncia acima devem ser demonstradas aqui em uma tabela com o cronograma destas atividades.>}

\subsection*{Li es Aprendidas}

\hl{<Citar as li es aprendidas durante o desenvolvimento das atividades e justificar a import ncia destas li es aprendidas.>}

\subsection*{Dificuldades Encontradas}

\hl{<Citar as dificuldades encontradas durante o desenvolvimento das atividades e mostrar se elas foram superadas e como foram superadas.>}

\subsection*{Oportunidades de Melhoria}

\hl{<Nada est  t o bom que n o possa melhorar. Pense como voc  pode melhorar, o que pode aprender para desempenhar suas atividades de forma mais satisfat ria ainda.>}


\addcontentsline{toc}{section}{Conclus o}
\section*{Conclus o}

\hl{<Fa a um encerramento de seu relat rio.>}
\end{document}